\documentclass[10pt, a4paper]{article}

\usepackage{setspace} % 줄간격 패키지
\setstretch{1.3}      % 1.6배 (160%) 설정
\usepackage{fontspec}
\usepackage{geometry}
% 여백: 상하좌우 1.2cm
\geometry{left=1.6cm, right=1.6cm, top=1.6cm, bottom=1.6cm}
\usepackage{titlesec}
\usepackage{enumitem}
\setlist{nolistsep} % 리스트 간격 제거
\usepackage{hyperref}
\usepackage{kotex}
\usepackage{array}
\usepackage{xcolor}
\usepackage{fontawesome}

\setmainfont{Roboto}
\setsansfont{Roboto}

% 색상 정의
\definecolor{bonusSteelBlue}{RGB}{70, 130, 180}
\definecolor{darkGray}{RGB}{60, 60, 60}
\definecolor{lightGray}{RGB}{150, 150, 150}

% 하이퍼링크 설정
\hypersetup{
    colorlinks=true,
    linkcolor=bonusSteelBlue,
    urlcolor=bonusSteelBlue,
}

% 섹션 스타일
\titleformat{\section}
{\large\bfseries\color{darkGray}}
{}
{0em}
{}[\titlerule]

% 커스텀 명령어
\newcommand{\jobtitle}[2]{\textbf{#1} \hfill \textit{#2}}
\newcommand{\dates}[1]{\hfill {\small\color{lightGray} #1}}

\begin{document}
\pagestyle{empty}

% huge --- HEADER START ---
\noindent
\begin{minipage}[t]{0.70\textwidth} % 왼쪽: 이름, 직무, 프로필
    {\huge\bfseries 두혁진}  \\[0.6em]
    {\Large\textcolor{bonusSteelBlue}{System Software Engineer}} 
    \vspace{0.8em}
    
    {\small 
    시스템의 기저 원리 탐구와 운영 자동화를 즐기는 엔지니어입니다. \\
    C/C++ 기반의 시스템 프로그래밍 역량과 Rust 백엔드 개발 경험으로 \\
    모빌리티 소프트웨어의 안정성에 기여하고 싶습니다.}
\end{minipage}%
\hfill
\begin{minipage}[t]{0.28\textwidth} % 오른쪽: 연락처
    {\huge\bfseries contact} \\
    {\small
    \faPhone \ 010-2805-0547 \\
    \faEnvelope \ \href{mailto:tyeirdoo06@gmail.com}{tyeirdoo06@gmail.com} \\
    \faGithub \ \href{https://github.com/hdoo42}{github.com/hdoo42}
    }
\end{minipage}

\vspace{1.0em}
% --- HEADER END ---


% --- 2단 레이아웃 본문 시작 ---
\noindent
\begin{minipage}[t]{0.34\textwidth} % 왼쪽 컬럼

    % 1. EDUCATION
    \section*{\huge{education}}
    42 Seoul \\ 
    \dates{2021.11 - 2023.12}
    
    \vspace{0.5em}
   
    % 한 문장 요약 (옵션 1 적용 예시)
    \small{C/C++ 기반의 견고한 시스템 프로그래밍(쉘·네트워크·동시성)부터 가상화 인프라 운영까지 컴퓨터 과학의 핵심 기초를 체화}\\
    \small{6대 coalition master(학생회장) 역임, 커뮤니티에 대한 기여를 인정받아 IITP원장상 수상}

    \vspace{1.0em}

    % 2. SKILLS
    \section*{\huge{ skills }}
    Programming \\
    \small{ Rust, C, C++, Bash, Python }
    
    \vspace{0.5em}
    Spoken Languages \\
    \small{ Japanese (Advanced Business) \\
    English (Daily / Simple Business) }    

    \vspace{0.5em}
    System \& OS \\
    \small{ Linux (Ubuntu), Shell Scripting }
    
    \vspace{0.5em}
    DevOps \& Tools \\
    \small{ Neovim, Git, Ansible, Docker, Prometheus, Grafana }

    \vspace{1.0em}

    % 3. ASK ME ABOUT
    \section*{\huge{ask me about}}
    Linux Kernel Troubleshooting
    \begin{itemize}[leftmargin=1em, label=\textbullet]
        \item Issue: Ubuntu 22.04 커널 업데이트 후 밝기 조절 불가 현상, Dell 기술지원이 해결 못한 문제
        \item Fix: \texttt{GRUB} 커널 파라미터 수정으로 해결
    \end{itemize}

\end{minipage}
\hfill % 좌우 컬럼 사이
\begin{minipage}[t]{0.62\textwidth} % 오른쪽 컬럼

    % 4. WORK EXPERIENCE
    \section*{\huge{work experience}}
    
    \jobtitle{(재)경산이노베이션아카데미}{System Software Engineer} \\
    \dates{2024.02 - 현재}
    \begin{itemize}[leftmargin=1em, label=\textbullet] 
        \item Operational Automation (GsRoot)
        \begin{itemize}
            \item Rust/SQLite 기반 슬랙봇 및 API 서버 개발로 워크스테이션 제어(세션 종료, 리셋) 자동화
            \item 관리자 수동 개입 최소화 및 24시간 무중단 운영 기여
        \end{itemize}
        \item System Integration (BLINK)
        \begin{itemize}
            \item 물리 보안 시스템(Biostar2)과 LDAP 인증 서버 간 미들웨어 개발
            \item 로그인-출입 권한 자동 동기화로 이기종 시스템 데이터 일관성 확보
        \end{itemize}
        \item Infrastructure Monitoring
        \begin{itemize}
            \item Prometheus, Grafana, Loki 도입으로 서버 및 클러스터 관제 시스템 구축
            \item 리소스 시각화로 장애 감지 시간 단축
        \end{itemize}
        \item Wiki Construction (codevalley.kr)
        \begin{itemize}
            \item 사내 기술 지식 공유를 위한 위키 구축 및 문서화 주도
        \end{itemize}
    \end{itemize}

    \vspace{0.8em}

    \jobtitle{(재)이노베이션 아카데미}{Intern} \\
    \dates{2023.10 - 2023.12}
    \begin{itemize}[leftmargin=1em, label=\textbullet]
        \item Ansible을 활용한 대규모 클러스터 프로비저닝 및 구성 관리 자동화
        \item 교육생 과제 명세 작성
    \end{itemize}

    \vspace{1.0em}

    % 5. KEY PROJECTS
    \section*{\huge{key projects}}
    \vspace{0.5em}
    \textbf{42 Region (Private Cloud Build)}
    \begin{itemize}[leftmargin=1em, label=\textbullet]
        \item OpenStack 기반 프라이빗 클라우드 구축 시도 중 네트워크 이슈(L1/L2 VLAN) 트러블슈팅 경험
        \item 인프라 네트워킹 및 가상화 기술에 대한 심도 있는 학습 계기 마련
    \end{itemize}

\end{minipage}

\end{document}
